\scanpage{274}
- \mspage{80} -
\margin{IX}
\margin{\figure{\begin{tikzpicture}\draw (0,0)--(.2,0);\draw[thick] (.5,-.2)--(.5,.15)--(.4,.13);\draw (.5,0)--(.35,-.18);\end{tikzpicture}}}
6.4. Notons que lorsque \(\ell\neq p,\) nous avons mis en évidence dans 4.3 a) une fonction à valeurs entières remarquable sur le groupe d'inertie \(I\) de \(S,\) d'ailleurs indépendante du choix de \(\ell,\) et dont la connaissance, dans le cas d'un corps résiduel parfait, suffit à reconstruire l'exposant conducteur de \(A\) (4.5), en prenant le produit scalaire, sur un quotient assez grand \(\Gamma=I/I'\) de \(I,\) de la fonction précédente et du conducteur d'Artin sur \(\Gamma\) (cf. (4.6.2)). Il peut être intéressant d'observer que cette fonction à valeurs entières sur \(I\) peut se réconstituer également par la connaissance de \(T_\ell(A_K)\) pour \(\ell=p\) - du moins (pour le moment) en inégales caractéristiques. Lorsque le corps résiduel \(k\) est fini, une observation analogue s'applique à la fonction à valeurs dans \(\mathbb Q,\) définie sur le sous-groupe \(\pi^!\) de \(\pi\) formé des éléments dont l'image dans \(\pi_0\) est de la forme \(f_q{}^n\) (\(n\in\mathbb Z\)), considérée dans 4.3 b). Indiquons le principe de la vérification de ces assertions. On utilise la filtration (6.3.1), et on est réduit à attacher une fonction convenable (sur \(I\) resp. sur \(\pi^!\)) à un \(p\)-pro-groupe de Barsoti-Tate \(Y\) sur \(K\) qui a ``essentiellement bonne réduction sur \(S\)'', au sens précisé dans 6.5. (\(Y\) sera un des trois quotients provenant de la filtration (6.3.1).) Dans le cas de 4.3 a), on se ramène au cas d'un corps résiduel parfait (ou même algébriquement clos, si on veut) par extension résiduelle convenable dans \(S\) (EGA 0\(_{\mathrm{III}}\) 10.3.1). Prenons alors une extension galoisienne finie \(K'\) de \(K\) telle que \(Y_{K'}\) ait bonne réduction sur \(S',\) d'où un groupe de Barsotti-Tate \(Y_{S'}\) sur \(S'\) et une opération de \(G\) sur \(Y_{S'}\) compatible avec son action sur \(S'.\) Si \(k'\) est le corps résiduel de \(S',\) \(G\)
\figure{\begin{tikzpicture}\draw[thick] (0,0) to[bend left=70] (.15,.1) to[bend left=70] (.25,0) to[bend left=70] (0,0);\draw (.9,0)--(.9,.2) (.8,.1)--(1,.1);\draw (2,0) to[bend left] (2.1,.2);\end{tikzpicture}}
\margin{398}
\margin{273}
\margin{398}
